\documentclass[10pt,a4paper]{amsart}
\usepackage[margin=19mm]{geometry}
\usepackage{amsmath,amssymb,mathtools}
\usepackage[scr=rsfs,cal=euler]{mathalfa}
\usepackage{xfrac}
\usepackage[unicode,hidelinks,bookmarks=false]{hyperref}
\newcommand{\ie}{i.e.\ }
\newcommand{\cf}{cf.\ }
\newcommand{\resp}{respectively\ }
\newcommand{\Subsection}{\subsection}
\providecommand{\nobreakdash}{\nobreak\hbox{-}}
\setlength{\emergencystretch}{2em}
\multlinegap0pt
\usepackage{bm}
\usepackage{bbm}
\usepackage{dsfont}
\usepackage{xspace}
\usepackage{cancel}
\usepackage{mathtools}
\usepackage{amsthm}
\makeatletter
\@ifundefined{lem}{\newtheorem{lem}{Lemma}}{}
\@ifundefined{thm}{\newtheorem{thm}{Theorem}}{}
\makeatother
\def \e {{\mathbf e}}
\def \o {{\bf 0}}
\newcommand{\cb}[1]{\left\lbrace #1\right\rbrace}
\newcommand{\ob}[1]{\left(#1\right)}
\title[$q$-Laplacian State Transfer on Graphs with Involutions]{$q$-Laplacian State Transfer on Graphs with Involutions}
\author{Swornalata Ojha, Hiranmoy Pal}
\date{Source: arXiv:2509.20749v1 (CC BY 4.0)}
\begin{document}
\maketitle

\noindent\emph{Source and licence.} $q$-Laplacian State Transfer on Graphs with Involutions. Version arXiv:2509.20749v1 (CC BY 4.0), distributed under the Creative Commons Attribution 4.0 International licence, as recorded in the arXiv licence metadata for this exact version and shown on the abstract page https://arxiv.org/abs/2509.20749v1. Full rights notice: licenses/arxiv-2509.20749-CC-BY-4.0.txt.\par\smallskip

\noindent\textit{Editorial scope.} Counted unit: the Thm whose source label is \texttt{T2} in the article (source0.tex lines 266--274, source label \texttt{T2}), delivered together with the article's own complete original proof of it (source0.tex lines 275--304, one \texttt{proof} environment of 30 lines). No mathematics is written, repaired or re-derived: statement and proof are the article's own text line for line. Required context retained uncounted: L2 (lines 244--255). Every definition, display and label of those lines is retained unchanged. Omitted from this package and not redistributed: 1--243, 256--265, 305--813. Nothing inside a retained excerpt is omitted or altered: every retained body is delivered exactly as the article's lines are written, so no omission receipt of any kind accompanies this package. Citations: the retained excerpts carry no citation command at all, so no bibliography entry of the article is transcribed and no reference is redistributed. Reference adaptation: none was needed and none was made; every reference inside the delivered unit resolves inside it, so no unresolved reference or citation is produced. Printed slips kept literal: no printed word, symbol, hypothesis or number of the counted statement or of its proof was altered. Layout and class: the article's own class is \texttt{article} with packages including etex, inputenc, amsmath, amssymb, amsthm, amsfonts, mathrsfs, systeme, bbm, commath, multicol, graphicx, tikz-network, authblk; the wrapper is the lane's pinned \texttt{amsart} editorial wrapper at 10pt on A4, which restates the theorem-like environments the retained text needs and adds the article's own macro definitions that the retained text needs (\texttt{\textbackslash cb}, \texttt{\textbackslash e}, \texttt{\textbackslash o}, \texttt{\textbackslash ob}), with extra wrapper packages (bm, bbm, dsfont, xspace, cancel, mathtools). The delivered numbering is the unit's own. Assets: no retained span contains a figure environment, a graphics-inclusion command, a PDF-page inclusion, a table or a TikZ picture; no image file is copied into this package, the package declares no asset dependency and no image file is redistributed with it. Licence: version arXiv:2509.20749v1 of this article is distributed under the Creative Commons Attribution 4.0 International licence, as recorded in the arXiv licence metadata for this exact version and shown on the abstract page https://arxiv.org/abs/2509.20749v1. A full rights notice is delivered at licenses/arxiv-2509.20749-CC-BY-4.0.txt. Characters: every retained excerpt is pure ASCII, so the delivered engine, XeLaTeX with its default Latin Modern text font, reports no missing character.
 \begin{lem}\label{L2}
 Let $G$ be a graph equipped with a non-trivial involution $\phi$, and $G'$ be a half graph induced by $\phi$. If $U_\mathscr{L}(t)$ is the transition matrix associated with the $q$-Laplacian matrix of $G,$ then for $u \in V(G')$,
 \begin{equation}\label{E3}
 U_\mathscr{L}(t)\ob{\e_u-\e_{\phi(u)}}=(I-P)\sum_{\mu_r \in \sigma(\mathscr{L}_-)}e^{it\mu_r}F_r\e_u,
 \end{equation}
 where P is the matrix of the involution $\phi$, and if $F_r'$ denotes the orthogonal projection of $\mathscr{L}_-$ corresponding to the eigenvalue $\mu_r,$ then the matrix $F_r$ is given by
 $$F_r=\frac{1}{2}\begin{bmatrix}
 F_r'&-F_r'&\o\\
 -F_r'& F_r'& \o\\
 \o& \o &\o
 \end{bmatrix}.$$
 \end{lem}

\begin{thm}\label{T2}
Let $G$ be a graph equipped with a non-trivial involution $\phi$, and $G'$ be a half graph induced by $\phi$. Suppose $M$ is the matrix whose columns are the vectors in $\mathcal{B}= \cb{\frac{1}{\sqrt{2} }\ob{\e_u-\e_{\phi(u)} }~|~  u \in V(G')},$ $\mathcal{C}=\cb{\frac{1}{\sqrt{2} }\ob{\e_u+\e_{\phi(u)} }~|~ u \in V(G')},$ and $\mathcal{D}=\cb{\e_u~|~ u \in S},$ respectively. If $\mathscr{L}$ is the $q$-Laplacian matrix associated with $G$, then for every $t \in \mathbb{R},$
\[M^T U_\mathscr{L}(t) M= \begin{bmatrix}	U_{\mathscr{L}_-}(t) & \o\\	\o & U_{\widetilde{\mathscr{L}_+}}(t)
\end{bmatrix},\]
where $\mathscr{L}',~\mathscr{L}_{-},~\mathscr{L}_{S}, ~A_\phi, ~A_S$ are as defined previously, and
\[\widetilde{\mathscr{L}_+}=\begin{bmatrix}	\mathscr{L}'+A_\phi&\sqrt{2} A_S\\
\sqrt{2}A_S^T & \mathscr{L}_S
\end{bmatrix}.\]
\end{thm}

\begin{proof}
Let $\e_u$ and $\check{\e}_u$ be the characteristic vectors of a vertex $u$ in $G$ and $G',$ respectively. Since $U_\mathscr{L}(t)$ is a polynomial in $\mathscr{L},$ the permutation matrix $P$ associated with the involution $\phi$ commutes with $U_\mathscr{L}(t).$ Since $\ob{\e_v-\e_{\phi(v)}}=\ob{I-P}\e_v$, the following holds for every $u,v\in V\ob{G'}$
\begin{eqnarray*}
 \frac{1}{2}\ob{\e_v-\e_{\phi(v)}}^T U_\mathscr{L}(t)\ob{\e_u-\e_{\phi(u)}} 
 %& = & \frac{1}{2}\e_v^T (I-P) U_\mathscr{L}(t) (I-P)\e_u\\
 &=& \e_v^TU_\mathscr{L}(t)\ob{\e_u-\e_{\phi(u)}}.
\end{eqnarray*}
It follows from Lemma \ref{L2} that
\begin{eqnarray*}
  \e_v^TU_\mathscr{L}(t)\ob{\e_u-\e_{\phi(u)}} &=& \e_v^T(I-P)\sum_{\mu_r \in \sigma(\mathscr{L}_-)}e^{it\mu_r}~F_r\e_u\\
  &=& \frac{1}{2}\sum_{\mu_r \in \sigma(\mathscr{L}_-)}e^{it\mu_r}~ \begin{bmatrix}
      \check{\e}_v^T& -\check{\e}_v^T& \o
  \end{bmatrix} \begin{bmatrix}
   F_r'\\
   -F_r'\\
   \o
  \end{bmatrix}\check{\e}_u\\
%  &=& \check{\e}_v^T\sum_{\mu_r \in \sigma(\mathscr{L}_-)}e^{it\mu_r}~ F_r' \check{\e}_u\\
  &=& \check{\e}_v^T U_{\mathscr{L}_-}(t) \check{\e}_u.
\end{eqnarray*}
Since $W = \text{span}(\mathcal{B})$ is an invariant subspace of $U_{\mathscr{L}}(t)$, the matrix representation of the restriction operator $U_{\mathscr{L}}(t)|_W$ relative to the basis $\mathcal{B}$ is the transition matrix $U_{\mathscr{L}_-}(t)$. Let $N$ be the matrix whose columns are the vectors in $\mathcal{C}$ and $\mathcal{D}$, respectively. Then
\[
N = \begin{bmatrix}
\frac{1}{\sqrt{2}} I_{|V(G')|} & \mathbf{0} \\
\frac{1}{\sqrt{2}} I_{|V(G')|} & \mathbf{0} \\
\mathbf{0} & I_{|S|}
\end{bmatrix}.
\]
It follows that $\mathscr{L}N=N\widetilde{\mathscr{L}}_+$, and therefore $U_\mathscr{L}(t)N=NU_{\widetilde{\mathscr{L}_+}}(t).$ Since the columns of $N$ span the orthogonal complement of $W$, we have the desired conclusion.
\end{proof}

\end{document}
